\section{Линейный оператор: основные понятия}

\begin{definition}
Отображение $\mathcal A\colon V\to W$ линейных пространств называется \emph{линейным оператором},
если
\[
  \mathcal A(x+y)=\mathcal Ax+\mathcal Ay,\qquad \mathcal A(\alpha x)=\alpha\,\mathcal Ax
\]
(аддитивность и однородность). Эквивалентно: $\mathcal A(\alpha x+\beta y)=\alpha\mathcal Ax+\beta\mathcal Ay$.
\end{definition}

\begin{property}\leavevmode
$\mathcal A0=0$; $\mathcal A(-x)=-\mathcal Ax$; $\mathcal A\bigl(\sum\alpha_i x_i\bigr)=\sum\alpha_i\mathcal Ax_i$.
\end{property}

\begin{example}
Поворот плоскости, дифференцирование $\mathcal D\colon P_n\to P_{n-1}$, умножение на матрицу
$x\mapsto Ax$ в $\R^n$.
\end{example}

\section{Ядро и образ. Теорема о ранге и дефекте}

\begin{definition}
\emph{Ядро} $\Ker\mathcal A=\{x\in V:\mathcal Ax=0\}$, \emph{образ}
$\Img\mathcal A=\{\mathcal Ax:x\in V\}\subseteq W$. \emph{Ранг} $\rang\mathcal A=\dim\Img\mathcal A$,
\emph{дефект} $\mathrm{def}\,\mathcal A=\dim\Ker\mathcal A$.
\end{definition}

\begin{theorem}
$\Ker\mathcal A$ — подпространство $V$, $\Img\mathcal A$ — подпространство $W$.
\end{theorem}

\begin{proof}
Если $x_1,x_2\in\Ker\mathcal A$, то $\mathcal A(\alpha x_1+\beta x_2)=\alpha\cdot0+\beta\cdot0=0$,
т.е.\ $\Ker\mathcal A$ замкнуто относительно линейных комбинаций. Если $y_i=\mathcal Ax_i\in\Img$,
то $\alpha y_1+\beta y_2=\mathcal A(\alpha x_1+\beta x_2)\in\Img$.
\end{proof}

\begin{theorem}[о ранге и дефекте]
$\mathrm{def}\,\mathcal A+\rang\mathcal A=\dim V$.
\end{theorem}

\begin{proof}
Возьмём базис $e_1,\dots,e_k$ ядра и дополним его до базиса $e_1,\dots,e_k,f_1,\dots,f_m$ всего $V$
($k+m=\dim V$). Покажем, что $\mathcal Af_1,\dots,\mathcal Af_m$ — базис образа. Они порождают
$\Img\mathcal A$: для $x=\sum a_ie_i+\sum b_jf_j$ имеем $\mathcal Ax=\sum b_j\mathcal Af_j$
(слагаемые с $e_i$ обнуляются). Независимы: если $\sum c_j\mathcal Af_j=0$, то
$\mathcal A(\sum c_jf_j)=0$, значит $\sum c_jf_j\in\Ker\mathcal A=\Span(e_i)$; но $f_j$ дополняют
базис ядра, поэтому все $c_j=0$. Итак $\rang\mathcal A=m$, и $k+m=\dim V$.
\end{proof}

\section{Матричное представление линейного оператора}

\begin{definition}
Пусть в $V$ выбран базис $\{e_j\}$, в $W$ — базис $\{f_i\}$. \emph{Матрица оператора}
$A=(a_{ij})$ задаётся разложениями образов базисных векторов: $\mathcal Ae_j=\sum_i a_{ij}f_i$
(столбец $j$ матрицы — координаты $\mathcal Ae_j$).
\end{definition}

\begin{theorem}
Если $X$ — столбец координат $x$, $Y$ — координат $\mathcal Ax$, то $Y=AX$. Отображение
$\mathcal A\mapsto A$ — изоморфизм пространств $\mathcal L(V,W)\cong\mathrm{Mat}_{m\times n}$
($n=\dim V$, $m=\dim W$).
\end{theorem}

\begin{proof}
$\mathcal Ax=\sum_j x_j\mathcal Ae_j=\sum_j x_j\sum_i a_{ij}f_i=\sum_i\bigl(\sum_j a_{ij}x_j\bigr)f_i$,
т.е.\ $Y_i=\sum_j a_{ij}x_j$ — это $Y=AX$. Соответствие сохраняет сложение и умножение на скаляр и
биективно (по матрице оператор восстанавливается на базисе), значит изоморфизм.
\end{proof}

\section{Ранг матрицы. Обратимость оператора}

\begin{theorem}
Ранг матрицы оператора равен рангу оператора и не зависит от выбора базисов. Оператор
$\mathcal A\colon V\to V$ обратим $\iff\det A\ne0\iff\rang A=\dim V$.
\end{theorem}

\begin{proof}
Столбцы $A$ — координаты $\mathcal Ae_j$, их линейная оболочка — образ; число независимых столбцов
$=\dim\Img\mathcal A=\rang\mathcal A$. Обратимость равносильна $\Ker\mathcal A=\{0\}$, т.е.\
$\mathrm{def}=0$, что по теореме о ранге даёт $\rang=\dim V$ (полный ранг, $\det A\ne0$).
\end{proof}

\section{Эндоморфизмы и алгебра операторов}

\begin{definition}
\emph{Эндоморфизм} — линейный оператор $\mathcal A\colon V\to V$. На $\mathcal L(V,V)$ заданы
операции: $(\mathcal A+\mathcal B)x=\mathcal Ax+\mathcal Bx$, $(\alpha\mathcal A)x=\alpha\mathcal Ax$,
композиция $(\mathcal A\mathcal B)x=\mathcal A(\mathcal Bx)$.
\end{definition}

\begin{remark}
$\mathcal L(V,V)\cong\mathrm{Mat}_{n\times n}$ как алгебры: композиции операторов отвечает
произведение матриц. \emph{Многочлен от оператора}
$p(\mathcal A)=c_k\mathcal A^k+\dots+c_1\mathcal A+c_0\mathcal I$.
\end{remark}

\section{Обратимые линейные операторы}

\begin{theorem}[критерий обратимости]
Для эндоморфизма конечномерного пространства инъективность, сюръективность и обратимость
равносильны: $\Ker\mathcal A=\{0\}\iff\Img\mathcal A=V\iff\exists\mathcal A^{-1}$.
\end{theorem}

\begin{proof}
По теореме о ранге $\dim\Ker+\dim\Img=\dim V$, поэтому $\Ker=\{0\}\iff\dim\Img=\dim V\iff\Img=V$.
Биективность даёт обратное отображение, линейность которого проверяется:
$\mathcal A^{-1}(\alpha y_1+\beta y_2)=\alpha\mathcal A^{-1}y_1+\beta\mathcal A^{-1}y_2$. Матрица
обратного оператора есть $A^{-1}$.
\end{proof}

\section{Смена базиса. Подобие матриц}

\begin{definition}
\emph{Матрица перехода} $S$ от базиса $\{e_i\}$ к $\{e'_j\}$: столбцы $S$ — координаты новых
базисных векторов в старом базисе, $e'_j=\sum_i s_{ij}e_i$. Координаты связаны как $X=SX'$. $S$
обратима.
\end{definition}

\begin{theorem}
При смене базиса матрица эндоморфизма преобразуется по закону подобия
\[
  A'=S^{-1}AS .
\]
След, определитель и ранг — инварианты подобия.
\end{theorem}

\begin{proof}
В старом базисе $Y=AX$; в новом $X=SX'$, $Y=SY'$. Тогда $SY'=ASX'$, откуда
$Y'=S^{-1}AS\,X'$, т.е.\ $A'=S^{-1}AS$. Инвариантность:
$\det A'=\det S^{-1}\det A\det S=\det A$; $\tr(S^{-1}AS)=\tr(ASS^{-1})=\tr A$ (циклическое свойство
следа); ранг сохраняется при умножении на обратимые матрицы.
\end{proof}
